\documentclass[12pt,a4paper]{article}
\usepackage[spanish,es-nodecimaldot]{babel}
\usepackage[utf8]{inputenc}
\usepackage[T1]{fontenc}
\usepackage{lmodern}
\usepackage[top=0.9cm,bottom=0.9cm,left=1.0cm,right=1.0cm]{geometry}
\usepackage{amsmath,amssymb}
\usepackage[table]{xcolor}
\usepackage{booktabs,array}
\usepackage{enumitem}
\usepackage{tikz}
\usetikzlibrary{arrows.meta}

\setlength{\parindent}{0pt}
\setlength{\parskip}{2.5pt}
\pagestyle{empty}

\newcommand{\cuerpo}{\fontsize{9.2}{11.1}\selectfont}

\AtBeginDocument{%
  \setlength{\abovedisplayskip}{6pt}%
  \setlength{\belowdisplayskip}{6pt}%
  \setlength{\abovedisplayshortskip}{3pt}%
  \setlength{\belowdisplayshortskip}{3pt}%
  \cuerpo
}

\begin{document}
\shorthandoff{<>}

\hfill{\bfseries\footnotesize Analisis de varianza}

\vspace{-4pt}
\begin{center}
{\fontsize{15}{18}\selectfont\bfseries Diseños completamente aleatorios}
\end{center}

La prueba ANOVA busca evaluar si las $k$ medias poblacionales $\mu_1,\mu_2,\ldots,\mu_k$ son todas iguales entre sí:

\begin{center}
\begin{tabular}{@{}l l@{}}
$H_0: \ \mu_1=\mu_2=\cdots=\mu_k$ & \hspace{0.6cm}(Todas las medias poblacionales son iguales) \\
$H_1: \ $ Al menos una media $\mu_i$ es distinta & \hspace{0.6cm}(Existe al menos un par $\mu_i \neq \mu_j$ para $i\neq j$) \\
\end{tabular}
\end{center}

\hspace{16pt}\begin{minipage}{0.94\linewidth}
Suponemos poblaciones normales con misma varianza, observaciones independientes. Podemos construir una matriz con k filas y n columnas

\vspace{6pt}
Primero considero el diseño completamente aleatorio con tamaños muestrales iguales. Tengo k tratamientos y cada uno tiene la misma cantidad n de observaciones. Las unidades experimentales se asignan aleatoriamente a cada tratamiento.
\end{minipage}

\begin{center}
{\fontsize{8.8}{10.6}\selectfont
\setlength{\tabcolsep}{5pt}
\begin{tabular}{l c c c c c c | c c}
\toprule
Tratamiento / Muestra & Obs. 1 & Obs. 2 & $\ldots$ & Obs. $j$ & $\ldots$ & Obs. $n$ & Total ($T_{i.}$) & Media ($\bar{y}_{i.}$) \\
\midrule
Muestra 1 & $y_{11}$ & $y_{12}$ & $\ldots$ & $y_{1j}$ & $\ldots$ & $y_{1n}$ & $T_{1.}$ & $\bar{y}_{1.}$ \\
Muestra 2 & $y_{21}$ & $y_{22}$ & $\ldots$ & $y_{2j}$ & $\ldots$ & $y_{2n}$ & $T_{2.}$ & $\bar{y}_{2.}$ \\
$\ldots$ & $\ldots$ & $\ldots$ & $\ldots$ & $\ldots$ & $\ldots$ & $\ldots$ & $\ldots$ & $\ldots$ \\
Muestra $i$ & $y_{i1}$ & $y_{i2}$ & $\ldots$ & $y_{ij}$ & $\ldots$ & $y_{in}$ & $T_{i.}$ & $\bar{y}_{i.}$ \\
$\ldots$ & $\ldots$ & $\ldots$ & $\ldots$ & $\ldots$ & $\ldots$ & $\ldots$ & $\ldots$ & $\ldots$ \\
Muestra $k$ & $y_{k1}$ & $y_{k2}$ & $\ldots$ & $y_{kj}$ & $\ldots$ & $y_{kn}$ & $T_{k.}$ & $\bar{y}_{k.}$ \\
\midrule
Total / Gran Media & & & & & & & $T_{..}$ & $\bar{y}_{..}$ \\
\bottomrule
\end{tabular}
}
\end{center}

\textbf{Efecto tratamiento ($\alpha$)} Interpretacion

\begin{minipage}[t]{0.60\linewidth}
\vspace{0pt}
Si \ $\bar{y}_{i.}$ \ (media muestra 1) tienen un valor mas grande que \ $\bar{y}_{..}$
\end{minipage}\hfill
\begin{minipage}[t]{0.38\linewidth}
\vspace{0pt}
{\small (media global o gran media) poniendo como ejemplo los detergente quiere decir q ese detergente es positivo o mejor}
\end{minipage}

Decir q las medias son iguales, los efectos serian 0.

\textbf{Para el calculo final del estadistico se realizan los sig calculos}

\textbf{1)} vamos a definir un valor constante que es la suma de los elementos de la matriz elevado al cuadrado, multiplicado por 1/k.n

\begin{center}
$\displaystyle C = \frac{1}{k\cdot n}\left(\sum_{i=1}^{k}\sum_{j=1}^{n} y_{ij}\right)^2$
\quad \textbf{o simplemente haciendo-{-}$>$} \quad
$\displaystyle C = \frac{(T_{..})^2}{k\cdot n}$
\hfill \textbf{C= termino de correccion}
\end{center}

\textbf{2)} vamos a calcular la suma de cuadrados total SST, tomando cada elemento de la matriz y elevando al cuadrado y sumando todo y restando el valor de ``C'' que calculamos recien. Es el primer miembro de la expresion.
\[
SST = \sum_{i=1}^{k}\sum_{j=1}^{n} \left(y_{ij}\right)^2 - C
\]

\textbf{3)} vamos a calcular una sumatoria llamada \textbf{suma de cuadrados de tratamientos}. Pertenece al ultimo termino de mi expresion. $T_i$ es la suma de todos los elementos de una fila de la matriz
\[
SS(Tr) = n\sum_{i=1}^{k} (\bar{y}_{i.}-\bar{y}_{..})^2
\qquad \textbf{o} \qquad
SS(Tr) = \frac{1}{n}\sum_{i=1}^{k} (T_{i.})^2 - C
\]

\textbf{4)} podemos calcular la suma de cuadrados de error SSE con el SST y SS(Tr) calculados anteriormente
\[
SSE = SST - SS(Tr)
\]

\textbf{5)} finalmente defino el estadistico F , para zona critica

\vspace{2pt}
\begin{center}
\begin{minipage}[c]{0.22\linewidth}
{\small
media cuad.\\ tratamiento(MS(Tr)) $\searrow$

\vspace{10pt}
media\\ cuadr.error\\ (MSE) $\nearrow$
}
\end{minipage}
\begin{minipage}[c]{0.26\linewidth}
\[
F = \frac{\dfrac{SS(Tr)}{k-1}}{\dfrac{SSE}{k\cdot(n-1)}}
\ = \
\frac{\hat{\sigma}_H^2}{\hat{\sigma}_W^2}
\]
\end{minipage}
\begin{minipage}[c]{0.20\linewidth}
{\small
$\longrightarrow$ n1=k-1 gl

\vspace{10pt}
$\longrightarrow$ n2=k(n-1) gl
}
\end{minipage}
\begin{minipage}[c]{0.30\linewidth}
{\small
$\hat{\sigma}_H^2 \longrightarrow$ Estima varianza a lo alto

\vspace{6pt}
$\hat{\sigma}_W^2 \longrightarrow$ Estima varianza a lo ancho

\vspace{6pt}
\hspace{14pt} k= nº tratamientos
}
\end{minipage}
\end{center}

\begin{minipage}[t]{0.26\linewidth}
\vspace{14pt}
{\small
\textbf{PARA k=4; n=12}

Error=total-tratamientos
\begin{align*}
&= (n.k-1)-(k-1) \\[-1pt]
&= nk-1-k+1 \\[-1pt]
&= nk-k \\[-1pt]
&= k(n-1)
\end{align*}
\vspace{-14pt}

$\longrightarrow$
}
\end{minipage}\hfill
\begin{minipage}[t]{0.72\linewidth}
\vspace{0pt}
{\fontsize{8.8}{10.6}\selectfont
\setlength{\tabcolsep}{3.5pt}
\begin{tabular}{|l|c|c|l|c|}
\hline
\textbf{Fuentes de} & \textbf{Grados de} & \textbf{Suma de} & \textbf{Media Cuadrada} & \textbf{F} \\
\textbf{Variación} & \textbf{Libertad} & \textbf{Cuadrados} & & \\
\hline
\textbf{Tratamientos} & k-1 & $SS(Tr)$ & MS(Tr)=SS(Tr)/(k-1) & MS(Tr)/MSE \\
\hline
\textbf{Error} & $k.(n-1)$ & $SSE$ & MSE=SSE/k.(n-1) & \\
\hline
\textbf{Total} & $n.k-1$ & $SST$ & & \\
\hline
\end{tabular}
}
\end{minipage}

\end{document}
